\chapter{Preliminares: Análisis funcional y medida}
\label{ch:preliminares}

\textcolor{red}{COSAS QUE HACER:}

\textcolor{red}{- CAMBIAR LOS AMBIENTES PARA QUE TENGAN EL NOMBRE ENTRE PARENTESIS (DE MANERA AUTOMATICA???) Y LUEGO UN SALTO DE LINEA}

\textcolor{red}{- FORMATO!!!}

\textcolor{red}{- PONER LEMA, TEOREMA, PROPOSICIÓN, etc CON LA PRIMERA LETRA EN MAYÚSCULA.}

\textcolor{red}{- PONER := EN LAS DEFINICIONES CUANDO TOQUE.}

\textcolor{red}{- CAMBIAR $||$ POR lVert.}

\textcolor{red}{- PONER EN MAYUSCULA CADA LETRA DEL TITULO DE LOS TEOREMAS, PROPOSICIONES, ETC (O NO, PERO CASARSE CON UNA DECISION). TAMBIEN PENSAR SI PONER Teorema O teorema CUANDO LO REFERENCIO... PENSAR ESTO TAMBIEN PARA LOS TITULOS DE LAS COSAS...}

\textcolor{red}{- PONER TODAS LAS SUCESIONES COMO $( \cdot )$ EN VEZ DE COMO $\{ \cdot \}$.}

\textcolor{red}{- HACER INTROS DE CAPITULOS.}

\textcolor{red}{- ASEGURARME QUE EN LAS DEFINICIONES SOLO HAYA UNA DEFINICION, NO VARIAS.}

\textcolor{red}{- DECIDIR SI CUANDO CITO UN TEOREMA/PROP/COROLARIO/DEF/etc LO HAGO CON EQREF O REF...}

\textcolor{red}{- PONER O CITAR SEMIGRUPO DUAL}

\textcolor{red}{UNA PEQUEÑA INTRODUCCIÓN}

\textcolor{red}{USAR $\top$ PARA TRANSPOSICIÓN}

\section{Análisis funcional}
\label{ch:preliminares:section:AnalisisFunc}

\subsection{Espacios de Hilbert}
\label{ch:preliminares:section:AnalisisFunc:subsec:EspHilbert}

\section{Teoría de la medida}
\label{ch:preliminares:section:TeoriaDeLaMedida}

\subsection{La integral de Bochner}
\label{ch:preliminares:section:TeoriaDeLaMedida:subsec:Bochner}